% Generated draft: requires mathematical and source review.
\documentclass[11pt]{article}
\usepackage[margin=25mm]{geometry}
\usepackage{fontspec}
\setmainfont{Latin Modern Roman}
\usepackage{amsmath,amssymb,mathtools,unicode-math}
\setmathfont{Latin Modern Math}
\usepackage{enumitem,needspace,xurl,hyperref}
\setlength{\parindent}{0pt}
\setlength{\parskip}{5pt}
\setlength{\emergencystretch}{3em}
\newcommand{\N}{\mathbb{N}}
\newcommand{\Z}{\mathbb{Z}}
\newcommand{\Q}{\mathbb{Q}}
\newcommand{\R}{\mathbb{R}}
\newcommand{\C}{\mathbb{C}}
\newcommand{\F}{\mathbb{F}}
\newcommand{\A}{\mathbb{A}}
\newcommand{\PP}{\mathbb{P}}
\newcommand{\E}{\mathbb{E}}
\newcommand{\eps}{\varepsilon}
\newcommand{\dd}{\,\mathrm{d}}
\begin{document}
\section*{UnsolvedMath ID: 10800003}
\textbf{Source problem code:} catalogue entry
\subsection{Definitions and mathematical statement}
The problem concerns a deformation of a singularity with a Lyashko--Looijenga map on a parameter ball \(B_\varepsilon\). In the customary isolated-critical-point setting, a regular parameter \(u\) has \(\mu\) distinct critical values, and the map records their unordered configuration:
\[\operatorname{LL}(u)=\{c_1(u),\ldots,c_\mu(u)\}.\]
Choose the reference fiber and vanishing paths defining cycles \(\delta_i,\delta_j\) for two of these values, and write \(\langle\delta_i,\delta_j\rangle\in\Z\) for their intersection number in the source's convention. The singularity, deformation, fiber dimension and choices are not reproduced in the supplied excerpt.

There are two questions. First, are there restrictions on collisions of critical values more refined than the intersection-number obstruction discussed in the source? Second, is the following sufficiency assertion true: for any two vanishing cycles with
\[\langle\delta_i,\delta_j\rangle\in\{-1,0,1\},\]
can the corresponding two critical values be made to collide by a collision in value space whose lift through the Lyashko--Looijenga submersion stays in \(B_\varepsilon\)? In a path formulation, one seeks a lift \(\gamma\) starting at the given regular parameter, with \(\operatorname{LL}(\gamma(t))\) approaching the intended pairwise collision as \(t\to1\), and with a collision parameter \(\gamma(1)\in B_\varepsilon\). The vanishing cycles in the premise must correspond to that collision via the chosen transport paths, not an unrelated choice of cycles.

The exact quantifier over collision paths is unstated: the excerpt does not decide whether a prescribed path must lift or whether some path may be chosen, whether other critical values are fixed, nor how \(\varepsilon\) is constrained. The path description therefore records the needed lifting content without choosing these missing conditions. The research summary reports that intersection number \(\pm2\) gives an established nonliftability obstruction in the source/classical literature; the new question asks for more refined obstructions and for sufficiency in the \(0,\pm1\) cases. That supplied negative report is not a new theorem proved here.
\subsection{Short English statement}
Determine finer obstructions to colliding critical values, and whether intersection number \(0\) or \(\pm1\) guarantees a collision lift within the prescribed deformation ball. The singularity and precise collision-path quantifiers are missing from the excerpt.
\subsection{Sources}
\begin{itemize}
\item Local UnsolvedMath registry, permanent ID \(10800003\), exact source code \texttt{catalogue entry}.
\item Victor Vassiliev, \emph{A Few Problems on Monodromy and Discriminants}, Problem 1C: \url{https://amj.math.stonybrook.edu/html-articles/Files-2015-2024/15-11/}. Recorded access 2026-07-29; html-structure extraction.
\item The research summary mentions classical literature for the \(\pm2\) obstruction but supplies no separate citation.
\end{itemize}
\end{document}
